\section{结论与不足}

\subsection{主要研究结论}

\subsubsection{福州市社区养老服务存在空间差异}

总结老年需求和养老服务供给的空间分布差异。

\subsubsection{高需求低供给区域具有明显识别价值}

总结哪些类型区域需要优先补位。

\subsubsection{代际共享服务能够部分缓解养老服务压力}

总结共享服务修正后的变化。

\subsubsection{MCLP 能为服务点布局提供量化依据}

总结新增 10、20、30 个点位的优化效果。

\subsection{研究贡献}

\subsubsection{方法贡献}

提出代际共享修正的养老服务可达性测度方法。

\subsubsection{实践贡献}

为福州市社区养老服务补位提供可执行方案。

\subsubsection{政策贡献}

为代际友好型社区建设提供量化依据。

\subsection{研究不足}

\subsubsection{老年人口网格估计存在误差}

由于缺少社区级真实老年人口，只能通过人口栅格和区县比例估计。

\subsubsection{部分服务设施容量数据可能缺失}

如床位数、收费标准、服务人数等字段不完整。

\subsubsection{模型结果仍需结合实地调研验证}

补位点位应结合土地、财政、社区意愿和实际运营条件进一步评估。

% ================================================================
